\section{Related Work}
\label{sec:related-work}

\subsection{The Shift from Heuristics to MCDM in Fire Emergency}
The evolution of fire emergency handling has moved from experience-based protocols to structured Multi-Criteria Decision-Making (MCDM) \cite{MDPI_MCDM_2025}. Early applications focused on the Analytic Hierarchy Process (AHP) for ranking urban fire susceptibility \cite{GIS_Site_Selection_2025}. However, recent literature emphasizes the reduction of subjective bias through hybrid models such as CRITIC-TOPSIS and Spherical Fuzzy Sets, which are better equipped to handle the "indeterminacy" inherent in Natech events \cite{SphericalFuzzy_2023, Neutrosophic_PERT_2025}.

\subsection{Strategic Resource Allocation and Fleet Management}
A core pillar of fire optimization research is Resource Allocation in Fire Departments (RAFD). Systematic reviews identify over 400 variables influencing fleet distribution, ranging from demographic shifts to temporal demand patterns \cite{Eslamzadeh2022}. Recent studies in heterogeneous fleet allocation have introduced lexicographic goals to balance human safety against equipment maintenance costs \cite{Fleet_Optimization_2025, Queueing_Fire_2023}. Despite these advancements in 2D spatial coverage, the transition to active, micro-level tactical deployment inside buildings remains under-researched \cite{SAFER_Grants_Analysis}.

\subsection{Indoor Monitoring: IoT, Digital Twins, and IRI}
The emergence of Building Information Modeling (BIM) and Digital Twins has transformed indoor fire monitoring \cite{Zhang2024}. By utilizing IoT sensor arrays, researchers have developed an Improved Risk Index (IRI) to provide real-time hazard visualization in multi-story residential buildings \cite{Wang2024, Abdullahi2025}. Knowledge-graph-based approaches now allow for the autonomous inference of fire progression laws by fusing textual reports with live sensor data \cite{Smith2025}. However, these systems often function as passive monitors rather than active tactical solvers \cite{TaylorFrancis2025}.

\subsection{Mathematical Programming: From Wildfire to Indoor MILP}
The most mathematically rigorous response models are often found in wildfire suppression literature. The MORA toolkit, for instance, utilizes lexicographic optimization to prioritize life-safety over property loss in expansive geographical landscapes \cite{Orphanoudakis2025, Initial_Attack_Success_2024}. Recent trends involve the use of Mixed-Integer Linear Programming (MILP) to handle resource scarcity under high uncertainty \cite{GrandaChico2025}. 



MORAv2 adapts these wildfire programming principles to 3D indoor node networks. By incorporating a dynamic ``smoke factor'' derived from scoping reviews of human behavior in fires \cite{Kinkel2025}, the model accounts for visibility degradation and toxic exposure \cite{ReyesNorambuena2025, Fisher2024}. This allows for the calculation of optimal rescue paths that are not only the fastest but also the safest, bridging the gap between wildfire-scale optimization and indoor tactical reality \cite{PTS_Algorithm_2025, Chen2025}.